\chapter{Énergie du vide, inflation et fluctuations primordiales}

\section{Énergie du vide et espace de de Sitter}

Dans un espace plat, l’équation de Friedmann s’écrit
\begin{equation}
\left(\frac{\dot A}{A}\right)^2=\frac{8\pi G}{3}\rho.
\end{equation}
Pour la matière, \(\rho\) décroît comme \(A^{-3}\). Pour le rayonnement, \(\rho\) décroît comme \(A^{-4}\). L’énergie du vide, elle, ne se dilue pas :
\begin{equation}
\rho=\rho_0=\text{constant}.
\end{equation}
Si l’Univers ne contient que cette énergie du vide, le membre de droite est constant. En notant cette constante \(H^2\),
\begin{equation}
H^2=\frac{8\pi G}{3}\rho_0,
\qquad
H=\sqrt{\frac{8\pi G}{3}\rho_0}.
\end{equation}
La constante de Hubble est alors véritablement constante :
\begin{equation}
\frac{\dot A}{A}=H.
\end{equation}
Ainsi,
\begin{equation}
\dot A=HA,
\end{equation}
et la solution est exponentielle :
\begin{equation}
A(t)=C e^{Ht}.
\end{equation}
La distance entre des galaxies lointaines et non liées augmente de façon exponentielle. C’est l’espace de de Sitter.

Dans ces coordonnées, la géométrie est donnée par
\begin{equation}
d\tau^2=dt^2-e^{2Ht}\left(dx^2+dy^2+dz^2\right).
\end{equation}
L’énergie du vide est ici positive. Dans un espace plat, elle doit l’être, car le membre de gauche de l’équation de Friedmann est non négatif. Une énergie du vide négative ne peut apparaître que si l’on tient compte des termes de courbure, ce qui n’est pas le cas étudié ici.

Einstein a introduit la constante cosmologique dans un autre but : équilibrer la matière ordinaire et maintenir l’Univers statique.

\section{L’Univers en accélération aujourd’hui}

L’espace de de Sitter représente aussi la tendance de l’Univers actuel à long terme. Auparavant, le facteur d’échelle suivait l’évolution propre à un univers dominé par la matière :
\begin{equation}
A(t)\propto t^{2/3}.
\end{equation}
Il y a quelques milliards d’années, la courbe s’est infléchie vers le haut et a commencé à croître plus rapidement. Elle n’est pas encore parfaitement exponentielle, car d’autres formes d’énergie sont toujours présentes, mais elle tend à le devenir.

Ici, l’accélération signifie que
\begin{equation}
\ddot A>0.
\end{equation}
C’est une caractéristique observée de l’expansion à l’époque tardive.

Considérons maintenant l’horizon. La loi de Hubble s’écrit
\begin{equation}
v=HD.
\end{equation}
À la distance où la vitesse de récession atteint celle de la lumière,
\begin{equation}
D=\frac{c}{H}.
\end{equation}
On peut également réécrire cette relation sous la forme \(HD=c\).\editorialnote{Il s’agit d’une reformulation algébrique, à des fins éditoriales, de la relation énoncée oralement \(D=c/H\). La ligne supérieure du tableau de droite, sur la figure~\ref{fig:lecture07-horizon}, n’est que partiellement lisible.}

\begin{figure}[ht]
\centering
\includegraphics[width=0.62\textwidth]{lecture_07_figure_02.png}
\caption{Relation de l’horizon cosmique \(D=c/H\) dans l’expansion de de Sitter.}
\label{fig:lecture07-horizon}
\lectureframe{00:30:37}
\end{figure}

Cette distance est de l’ordre de dix à quinze milliards d’années-lumière, une valeur numériquement proche de l’âge de l’Univers, une fois celui-ci converti en distance au moyen de \(c\).

\begin{classroomqa}
\lecturetimestamp{00:14:10}
\audiencequestion{L’égalité approximative entre l’échelle de l’horizon et l’âge de l’Univers ne vient-elle pas simplement du fait que la tangente passe par l’origine ?}
\lecturerresponse{Ce serait rigoureusement exact pour \(A(t)\propto t\). Mais ce n’est pas vrai en général. Dans l’espace de de Sitter, \(H\) reste fixe tandis que l’âge continue d’augmenter ; l’égalité approximative actuelle n’est donc pas une identité mathématique.}
\end{classroomqa}

\section{Ce que signifie l’énergie du vide}

L’énergie du vide est décrite ici dans le langage de la mécanique quantique ordinaire. Pour un oscillateur harmonique, l’énergie s’écrit schématiquement
\begin{equation}
E\sim p^2+x^2.
\end{equation}
Classiquement, l’état fondamental correspond au point d’équilibre,
\begin{equation}
x=0,\qquad p=0,\qquad E=0.
\end{equation}
La mécanique quantique ne permet pas que \(x\) et \(p\) soient tous deux exactement fixés à zéro. Le principe d’incertitude impose des fluctuations, de sorte que l’énergie de l’état fondamental est positive. Pour l’oscillateur, on obtient
\begin{equation}
E_0=\frac{1}{2}\hbar\omega.
\end{equation}

On étend ensuite cette idée aux champs. L’Univers est rempli d’oscillateurs harmoniques : les oscillations du champ électromagnétique, ainsi que celles de tous les autres champs. Chaque champ possède une énergie de point zéro. Pour le champ électromagnétique, la densité d’énergie s’écrit schématiquement
\begin{equation}
u_{\mathrm{EM}}\sim E^2+B^2.
\end{equation}
Même en l’absence de photons réels, des oscillations irréductibles de point zéro subsistent. C’est une explication de l’énergie du vide. La seule propriété nécessaire ici est que cette énergie est présente et ne se dilue pas à mesure que l’Univers se dilate.

\begin{classroomqa}
\lecturetimestamp{00:16:15}
\audiencequestion{Qu’est-ce que cette énergie du vide, et pourquoi resterait-il quelque chose si l’on retirait la matière et le rayonnement ?}
\lecturerresponse{Parce que le vide est l’état fondamental des champs quantiques, et que l’énergie de l’état fondamental d’un système quantique n’est pas nécessairement nulle. Les oscillateurs harmoniques ont une énergie fondamentale positive, et les champs se comportent comme des ensembles d’oscillateurs. Ainsi, même en l’absence de particules réelles, une énergie de point zéro subsiste.}
\end{classroomqa}

Il ne s’agit pas du rayonnement de fond à trois degrés. Le fond diffus cosmologique est un rayonnement ordinaire. Sa densité d’énergie diminue parce que les photons se diluent et que chacun d’eux subit un décalage vers le rouge. L’énergie du vide est l’énergie résiduelle de point zéro qui demeure une fois les photons réels retirés.

\section{Ce que signifie \(H^{-1}\)}

Revenons maintenant à la relation entre l’inverse de la constante de Hubble et l’âge de l’Univers. Si
\begin{equation}
A(t)\propto t,
\end{equation}
alors
\begin{equation}
H=\frac{\dot A}{A}=\frac{1}{t}.
\end{equation}
Plus généralement, si
\begin{equation}
A(t)\propto t^p,
\end{equation}
alors
\begin{equation}
H=\frac{p}{t}.
\end{equation}
Dans le cas dominé par la matière, \(p=2/3\),
\begin{equation}
H=\frac{2}{3t}.
\end{equation}
Ainsi, dans les modèles cosmologiques en loi de puissance, \(H^{-1}\) est du même ordre que l’âge, sans lui être exactement égal.

L’espace de de Sitter est différent. Dans ce cas,
\begin{equation}
A(t)=C e^{\gamma t},
\end{equation}
et donc
\begin{equation}
H=\frac{\dot A}{A}=\gamma,
\end{equation}
une véritable constante. Ici, \(H^{-1}\) ne représente pas l’âge de l’Univers. C’est la constante de temps de la croissance exponentielle : le temps au bout duquel le facteur d’échelle est multiplié par \(e\). Dans l’exposé oral, il est dit à plusieurs reprises ``double'', puis cette formulation est corrigée ; l’énoncé précis est une multiplication par \(e\). Aujourd’hui, cette constante de temps est de l’ordre de dix milliards d’années.

\begin{classroomqa}
\lecturetimestamp{00:23:46}
\audiencequestion{Pourriez-vous revenir sur le lien entre l’âge de l’Univers et l’inverse de la constante de Hubble ?}
\lecturerresponse{Pour les lois de puissance, on obtient \(H=p/t\), donc \(H^{-1}\) suit l’âge à un facteur d’ordre un près. Mais dans l’espace de de Sitter, \(H\) est constante tandis que l’âge continue d’augmenter ; \(H^{-1}\) est donc seulement la constante de temps de l’expansion exponentielle, et non l’âge.}
\end{classroomqa}

\section{Énergie et horizon cosmique}

Si la densité d’énergie du vide reste constante tandis que l’Univers se dilate, on peut naturellement se demander si l’énergie totale du vide augmente. La réponse donnée ici est nuancée. En relativité générale, le bilan énergétique comprend une composante gravitationnelle ; la notion d’énergie totale est donc plus subtile qu’en mécanique élémentaire. Le point pratique est celui qui concerne l’horizon.

Dans l’espace de de Sitter exact, l’horizon cosmique a une taille fixe,
\begin{equation}
D=\frac{c}{H}
=\frac{c}{\sqrt{8\pi G\rho_0/3}}.
\end{equation}
Le volume à l’intérieur de cet horizon est donc fixe. L’énergie du vide qu’il contient vaut
\begin{equation}
E_{\mathrm{hor}}=\frac{4\pi}{3}D^3\rho_0.
\end{equation}
\editorialnote{Cela résume le calcul présenté oralement de 00:31:06--00:32:23 : prendre la taille de l’horizon, calculer le volume de la sphère, puis multiplier par \(\rho_0\).}
Comme \(D\) et \(\rho_0\) sont tous deux constants, cette quantité ne varie pas au cours du temps. Pour la portion de l’Univers qu’un observateur donné pourra jamais voir, l’énergie du vide accessible reste constante dans l’espace de de Sitter exact.

Des objets peuvent aussi franchir l’horizon cosmique, à peu près comme ils franchissent l’horizon d’un trou noir. Une fois de l’autre côté, ils ne peuvent plus entrer en contact causal avec nous.

\begin{classroomqa}
\lecturetimestamp{00:29:27}
\audiencequestion{Si la densité d’énergie du vide reste constante tandis que l’Univers se dilate, d’où vient l’énergie supplémentaire ?}
\lecturerresponse{Le cours répond en termes de relativité générale, et non en imaginant une boîte newtonienne. Dans l’espace de de Sitter exact, la taille de l’horizon reste fixe ; l’énergie du vide dans la région qui pourra jamais être observée reste donc fixe elle aussi.}
\end{classroomqa}

\section{Pourquoi l’inflation a été introduite}

On transpose ensuite le même mécanisme de de Sitter dans un passé lointain, mais avec une constante de temps bien plus courte, de l’ordre de
\begin{equation}
H^{-1}\sim 10^{-32}\,\mathrm{s}.
\end{equation}
Telle est l’idée de l’inflation : avant toute période directement observable, l’Univers a connu une phase d’expansion exponentielle.

La motivation principale tient à la platitude et à l’homogénéité de l’Univers. Les tranches spatiales paraissent très plates. À grande échelle, l’Univers est également très homogène. Plus frappant encore, la surface de dernière diffusion était déjà homogène ; cette homogénéisation ne peut donc pas s’être produite uniquement au cours de l’histoire visible ultérieure.

Un ballon dégonflé est plissé. Gonflez-le suffisamment et toute portion visible devient lisse. L’analogie du raisin sec, ou du pruneau, illustre la même idée : un objet plissé, étiré d’un facteur immense, paraît lisse sur toute portion visible. L’inflation est censée avoir produit exactement cet effet.

Cela répond aussi au problème de l’horizon sous la forme présentée ici. Des régions très éloignées les unes des autres sur la surface de dernière diffusion sont toutes deux visibles pour nous, mais la lumière n’a pas eu le temps de passer de l’une à l’autre au cours de l’histoire visible ultérieure. Leur densité et leur température presque uniformes suggèrent donc un passé durant lequel une région plus petite a été étirée de façon considérable.

L’inflation a un effet ambivalent. Elle efface la sensibilité aux conditions initiales et, pour la même raison, fait disparaître une grande partie des informations sur les débuts de l’Univers.

La question de la platitude se pose ici sous une forme plus précise. Du point de vue observationnel, le paramètre de courbure \(k\) semble nul, à un pour cent près, voire mieux. La géométrie paraît euclidienne dans la région visible. L’inflation l’explique en rendant le rayon de courbure bien plus grand que la région que nous pouvons observer.

\begin{classroomqa}
\lecturetimestamp{00:40:05}
\audiencequestion{Suppose-t-on que l’Univers était à l’équilibre avant l’inflation ?}
\lecturerresponse{Non. L’idée de l’inflation est presque à l’opposé. Quelle que soit la structure présente auparavant, elle est étirée à un tel point que la région visible perd la mémoire des détails antérieurs.}
\end{classroomqa}

\section{Un champ scalaire et son potentiel}

Pour expliquer comment l’énergie du vide effective a pu être bien plus élevée dans l’Univers primordial qu’elle ne l’est aujourd’hui, introduisons un champ scalaire. Contrairement aux champs électrique et magnétique, qui sont des vecteurs, un champ scalaire n’a pas de direction. Il possède simplement une valeur en chaque point :
\begin{equation}
\phi=\phi(x,y,z,t).
\end{equation}
Ce champ est hypothétique. Il n’est identifié à aucun champ connu dans la nature, et ses quanta ne sont identifiés à aucune particule connue.

La partie de son énergie qui nous intéresse ici est une énergie potentielle,
\begin{equation}
V(\phi).
\end{equation}
Par ``énergie potentielle'', on entend une énergie associée à la valeur du champ lui-même, et non à la dérivée temporelle du champ. Différentes valeurs de \(\phi\) peuvent correspondre à différentes énergies.

\begin{figure}[ht]
\centering
\includegraphics[width=0.68\textwidth]{lecture_07_figure_03.png}
\caption{Potentiel scalaire qualitatif \(V(\phi)\) montrant que différentes valeurs du champ correspondent à différentes énergies potentielles.}
\label{fig:lecture07-generic-potential}
\lectureframe{00:47:08}
\end{figure}

\begin{classroomqa}
\lecturetimestamp{00:46:34}
\audiencequestion{Que représente l’axe horizontal dans le dessin \(V(\phi)\) ?}
\lecturerresponse{L’axe horizontal représente la valeur du champ \(\phi\), et l’axe vertical, l’énergie potentielle \(V\).}
\end{classroomqa}

Si le champ est constant dans l’espace et dans le temps, ou presque constant dans la région considérée, alors \(V(\phi)\) se comporte exactement comme l’énergie du vide. Elle ne se dilue pas. On reprend donc les formules précédentes de de Sitter en remplaçant \(\rho_0\) par \(V(\phi)\). En particulier,
\begin{equation}
H=\sqrt{\frac{8\pi G}{3}V(\phi)}.
\end{equation}

Si \(\phi\) varie d’un endroit à l’autre, les différentes régions subissent une inflation avec des constantes de temps différentes. Mais si l’étirement est suffisamment important, toute portion visible devient presque homogène.

\begin{classroomqa}
\lecturetimestamp{01:05:11}
\audiencequestion{\(\phi\) joue-t-il ici le rôle d’un paramètre de temps ?}
\lecturerresponse{Dans cette discussion, oui, en ce sens limité que, pendant que le champ descend le potentiel, il est une fonction monotone du temps.}
\end{classroomqa}

\section{Un plateau, une falaise et le réchauffement}

Le potentiel nécessaire à ce scénario est très plat, mais pas parfaitement ; il atteint ensuite le bord d’une falaise, chute brutalement, sans tout à fait atteindre zéro, puis remonte.

\begin{figure}[ht]
\centering
\includegraphics[width=0.68\textwidth]{lecture_07_figure_04.png}
\caption{Potentiel inflationnaire avec un plateau élevé, une chute abrupte et un minimum plus bas.}
\label{fig:lecture07-inflationary-potential}
\lectureframe{01:02:30}
\end{figure}

Cette forme est très artificielle, mais elle est introduite parce que le scénario cosmologique semble exiger quelque chose de semblable.

Le champ commence en haut du plateau. Comme la pente est faible et que la viscosité effective associée à l’expansion est importante, \(\phi\) ne varie que très lentement. L’image visée est plutôt celle d’un rocher qui tombe lentement dans du miel que celle d’un objet qui glisse sur une table. Plus l’Univers se dilate rapidement, plus cette viscosité augmente.

\begin{classroomqa}
\lecturetimestamp{00:59:07}
\audiencequestion{Qu’est-ce qui fait évoluer \(\phi\) ?}
\lecturerresponse{La pente du potentiel le pousse. Cette pente est faible et la viscosité cosmologique est importante ; le champ dérive donc lentement plutôt que rapidement.}
\end{classroomqa}

Tandis que \(\phi\) avance lentement sur le plateau, \(V(\phi)\) varie très peu. Pour cette raison, l’Univers connaît une inflation presque comme si l’énergie du vide était constante. La constante de temps durant cette phase est de l’ordre de
\begin{equation}
H^{-1}\sim 10^{-32}\,\mathrm{s}.
\end{equation}
Une estimation orale approximative situe la hauteur du plateau près de \(10^{-14}\) en unités de Planck.\editorialnote{Il s’agit de l’estimation orale ``la plus plausible'', donnée à 01:14:38--01:14:47, et non d’une valeur déduite dans le cours.}

Combien de temps l’inflation a-t-elle duré ? On ne le sait pas, car ni la largeur ni la faible pente du plateau ne sont connues. Ce que l’on sait, c’est que l’Univers a été multiplié au moins par \(e\), environ cinquante fois, et que soixante est le nombre habituellement cité.\editorialnote{Reformulation éditoriale de la discussion orale sur le nombre d’e-foldings à 01:14:05--01:19:27.}Il s’agit d’e-foldings : chaque e-fold correspond à une multiplication par \(e\).

Lorsque le champ atteint le bord, l’énergie potentielle varie rapidement. Le champ chute ; son énergie potentielle est d’abord convertie en énergie cinétique, puis en chaleur et en particules, et l’Univers devient chaud. Le terme ``réchauffement'' est légèrement trompeur dans cette description, car toute énergie thermique antérieure aurait déjà été diluée pendant l’inflation.

Le minimum inférieur n’est pas à zéro. La petite valeur qui y subsiste est associée à la lente expansion exponentielle actuelle.

\begin{classroomqa}
\lecturetimestamp{01:19:48}
\audiencequestion{Le plateau pourrait-il être parfaitement plat, les perturbations faisant d’une manière ou d’une autre basculer le champ au bord ?}
\lecturerresponse{Non. Il ne peut pas être parfaitement plat. Il doit avoir une certaine inclinaison. S’il était trop plat, l’Univers réchauffé serait beaucoup trop inhomogène.}
\end{classroomqa}

``Big Bang'' est employé ici dans deux sens proches. Il peut désigner l’événement de réchauffement lui-même. Il peut aussi désigner la première étape directement visible, à savoir la surface de dernière diffusion. Le point essentiel demeure que tout ce qui est effectivement observé se situe bien en dessous du plateau inflationnaire.

\begin{classroomqa}
\lecturetimestamp{01:45:42}
\audiencequestion{Le réchauffement lui-même correspond-il au Big Bang ?}
\lecturerresponse{Si vous voulez, oui. Mais à ce moment-là, l’Univers était encore extrêmement chaud et opaque. La surface de dernière diffusion apparaît bien plus tard, après que l’Univers a suffisamment refroidi pour devenir transparent.}
\end{classroomqa}

\section{Fluctuations quantiques et germes des galaxies}

L’inflation rend l’Univers visible très homogène, mais pas parfaitement. En l’absence de fluctuations quantiques, le champ deviendrait complètement lisse. La mécanique quantique fait qu’il varie encore légèrement d’un endroit à l’autre.

Le schéma central représente l’espace ordinaire sur l’axe vertical et la valeur du champ \(\phi\) sur l’axe horizontal. Si le champ était parfaitement homogène, chaque point de l’espace aurait la même valeur de \(\phi\) à un instant donné, et tous les points évolueraient ensemble vers le bord. Toutes les régions basculeraient alors au-dessus de la falaise au même instant. Les fluctuations quantiques rompent cette parfaite synchronisation.

Deux images illustrent cette idée. La première est une petite bosse dans un champ par ailleurs homogène. La seconde est une rangée de soldats marchant bras dessus bras dessous vers une falaise. S’ils sont parfaitement alignés, ils tombent tous ensemble. Si certains sont légèrement en avance et d’autres légèrement en retard, ils basculent à des moments différents.

L’effet important intervient alors. Une région qui bascule plus tôt se réchauffe plus tôt. Quand une région plus tardive bascule à son tour, la première a déjà eu le temps de se dilater et de refroidir quelque peu. Cette différence engendre des variations de densité d’énergie après le réchauffement.

\begin{classroomqa}
\lecturetimestamp{01:34:00}
\audiencequestion{Les régions réchauffées plus tard sont-elles plus chaudes ou plus froides que les premières ?}
\lecturerresponse{Les premières ont déjà eu le temps de se dilater et de refroidir. Les régions plus tardives sont donc plus chaudes au moment où elles basculent, et cette différence contribue à l’inhomogénéité de densité.}
\end{classroomqa}

L’effet observé est faible :
\begin{equation}
\frac{\delta\rho}{\rho}\sim 10^{-5},
\end{equation}
environ un cent millième. Il s’agit de fluctuations primordiales de densité, et non d’énergie noire. Cette valeur n’est pas déduite ici de la seule esquisse qualitative ; elle sert à contraindre la hauteur et la platitude du plateau.

Elles sont également nécessaires. Sans elles, il n’y aurait pas de galaxies. Après le réchauffement, la gravité amplifie le motif initial. Les régions surdenses attirent la matière des régions sous-denses. Les régions surdenses deviennent encore plus denses, les régions sous-denses encore moins denses, et il en résulte la structure ultérieure des galaxies et des étoiles.

Il ne s’agit pas seulement d’un scénario qualitatif. On peut calculer les fluctuations quantiques d’un champ dans un Univers en expansion, les suivre à travers le réchauffement et comparer le résultat au motif des fluctuations reconstitué à partir de la distribution des galaxies. L’accord est très bon.

Cette discussion conduit aussi à une remarque sur l’entropie. L’agglomération gravitationnelle paraît contre-intuitive, mais elle ne viole pas le second principe. La gravité a un comportement particulier en mécanique statistique : elle tend à faire s’agglomérer les objets, et cette agglomération augmente l’entropie au lieu de la diminuer.

\begin{classroomqa}
\lecturetimestamp{01:40:24}
\audiencequestion{Quel rôle joue le fond diffus cosmologique dans tout cela, étant donné que \(\phi\) n’est pas un champ électromagnétique ?}
\lecturerresponse{Le motif primordial de densité s’imprime sur la surface de dernière diffusion. Le réchauffement convertit l’énergie du champ en particules ordinaires et en photons, puis le fond diffus cosmologique conserve la trace de ces fluctuations primordiales.}
\end{classroomqa}

\section{Ondes acoustiques et étape suivante}

À la fin, la discussion aborde brièvement les ondes acoustiques. Elles précèdent l’agglomération gravitationnelle ultérieure. Les fluctuations primordiales évoluent en oscillations acoustiques, qui laissent dans le ciel des concentrations de matière d’une taille physique connue.

C’est important, car une taille physique connue à une distance connue fournit une règle de mesure. En mesurant l’angle qu’elle sous-tend sur la surface de dernière diffusion, on peut déterminer, à l’aide de la géométrie, si l’espace est plat, à courbure positive ou à courbure négative. Les détails sont reportés.

\begin{classroomqa}
\lecturetimestamp{01:47:19}
\audiencequestion{Qu’en est-il des ondes acoustiques ?}
\lecturerresponse{Elles fournissent une règle de mesure dans le ciel. Leur taille physique connue, associée à leur taille angulaire observée sur la surface de dernière diffusion, permet de mesurer la géométrie de l’espace. Les détails sont reportés au prochain cours.}
\end{classroomqa}